\section*{2 Preliminaries}
\subsection*{2.1 Basic Notions}
We refer to the book by Arora and Barak~\cite{ref9} for definitions of standard complexity classes.
Given a circuit class $\mathcal C$, we use $\mathcal C[s(n)]$ to refer to the class of functions computed by $\mathcal C$-circuits of size $s(n)$.
Given a language $L\subseteq\{0,1\}^*$, $\operatorname{co}\text{-}L$ denotes the complement of $L$. Given language $L$ and $n\in\mathbb N$, $L_n$ is used to refer to $L\cap\{0,1\}^n$.
Given a set $S\subseteq\mathbb N$ and languages $L$ and $L'$, we say $L$ agrees with $L'$ on $S$ if for all $n\in S$, $L_n=L'_n$.
A time-constructible function $T:\mathbb N\to\mathbb N$ is a function such that there is a Turing machine transducer $M$, which for each $n$, on input $1^n$ halts with output $T(n)$ within $O(T(n))$ steps.
We need a generalization of the notion of robustly-often simulation from~\cite{ref11}. Given $g:\mathbb N\to\mathbb N$, a set $S\subseteq\mathbb N$ is called $g$-robust if for each constant $k$, there is $m\in\mathbb N$ such that for all $n$ such that $m\leq n\leq g(m)^k$, $n\in S$. Note that any $g$-robust set is also $\poly(g)$-robust.
Given a function $g$, a language $L$ and a complexity class $\mathcal C$, we say $L$ is $g$-robustly often in $\mathcal C$ if there is $L'\in\mathcal C$ for which there is a $g$-robust set $S$ such that $L$ agrees with $L'$ on $S$. Given function $g$ and complexity classes $\mathcal C$ and $\mathcal C'$, we say that $\mathcal C$ is $g$-robustly often in $\mathcal C'$ if for all languages $L\in\mathcal C$, $L$ is $g$-robustly often in $\mathcal C'$. We use the term ``robustly often'' as a synonym for $g$-robustly often with $g$ the identity function.
The proposition below is a parameterized version of Theorem12 in~\cite{ref11}.
\setcounter{nttheorem}{2}
\begin{ntproposition}[Author Proposition3]
Let $T$ be a time-constructible function such that $T(n)\geq n$ for all $n$. If SAT is $T(\poly(n))$-robustly often in $\DTIME(T)$, then $\Sigma_2\mathrm P$ is robustly often in $\DTIME(T(\poly(T(\poly(n)))))$.
\end{ntproposition}
The following is a simple application of standard diagonalization~\cite{ref13}.
\begin{ntproposition}[Author Proposition4]
Let $T$ be a time-construtible function such that $T=o(2^n)$. Then $\mathrm E$ is not robustly often in $\DTIME(T)$.
\end{ntproposition}
\subsection*{2.2 Resource-Bounded Kolmogorov Complexity}
We study various notions of resource-bounded Kolmogorov complexity, and associated decision problems.
Fix a universal Turing machine $U$. The $Kt$ complexity of a string is defined as follows:
\[Kt(x)=\min\{|p|+\log(t)\mid U(p)\text{ halts and outputs }x\text{ within }t\text{ steps}\}.\]
In the MKtP problem, the instance is a string $x$ together with a parameter $s$, and the question is whether $Kt(x)\leq s$.
It is known that MKtP is complete for $\mathrm E$ under polynomial-size truth-table reductions~\cite{ref2}, but completeness under uniform polynomial-time reductions is unknown.
In the MCSP problem, the instance is a string $x$ together with a parameter $s$, and the question is whether $x$ is the truth table of a Boolean function with circuit complexity at most $s$.
MCSP is easily seen to be in NP, but it is unknown whether MCSP is NP-hard. A brute-force search strategy of trying all circuits $C$ of size at most $s$ and checking if any of them computes the function with truth table $x$ can easily be implemented to run in time $\poly(|x|)s^{O(s)}$.
\subsection*{2.3 Reductions}
We will work with constrained notions of polynomial-time Turing reduction.
\begin{ntdefinition}[Author Definition5]
A polynomial-time Turing reduction from a language $L$ to MCSP or MKtP is called parametric honest if there is a constant $\epsilon>0$ such that for every $x\in\{0,1\}^*$, every query made by the reduction on $x$ has parameter bounded below by $|x|^\epsilon$.
\end{ntdefinition}
Note that parametric honesty is not implied by the standard notion of honesty for polynomial-time reductions, which only requires the output length to be bounded below by some polynomial in the input length. Implicit in the definition is that the lower bound on parameters of queries is only required to hold when previous queries are answered correctly; when queries are answered wrongly, ``all bets are off''.
\begin{ntdefinition}[Author Definition6]
A polynomial-time Turing reduction from a language $L$ to MCSP or MKtP is called natural if for any input $x\in\{0,1\}^*$, the parameter of any query made by the reduction on $x$ depends only on $|x|$.
\end{ntdefinition}
This notion of naturalness generalizes the notion used by Kabanets and Cai for $m$-reductions to Turing reductions. In fact, it is slightly weaker than their notion when restricted to $m$-reductions, since it does not constrain the length of the output, merely the parameter. As with the definition of parametric honesty, the condition on parameters is only required to hold when previous queries are answered correctly.
Naturalness is incomparable to parametric honesty--parametric honesty allows the parameter to vary but restricts its range of variation, while naturalness allows the parameter to have a large range of variation but insists that it is the same for all queries made on any input of a given length.
